\section{Discussion}

\subsection{What Should Be Reported}

Our evidence supports three reporting recommendations for repo-grounded agent
benchmarks.
\begin{enumerate}[leftmargin=1.25em]
\item Report the benchmark's official \processmetric and \resultmetric fields
separately whenever both are available, so readers can distinguish
\processmetric$=$False from \processmetric$=$True,
\resultmetric$=$False outcomes.
\item Report family or domain slices in addition to one global total.
\item Include representative official \processmetric$=$True,
\resultmetric$=$False comments for structured extraction tasks, not only
aggregate pass counts.
\end{enumerate}

These recommendations are intentionally conservative. We do not ask benchmark
authors to redesign their official metric. We ask evaluators to stop hiding
existing official structure behind one number.

\subsection{Limitations}

The paper has clear limits. First, the main evidence is local-run rather than a
synchronized scrape of the public \gittaskbench leaderboard. Second, broader
frontier API coverage currently exists only on the two-task shared
\texttt{Trafilatura} slice, where we keep the latest local run per model-task
pair rather than a synchronized same-day rerun. Third, some
\processmetric$=$False mass may still reflect local harness brittleness rather
than pure model weakness. Some \processmetric$=$True,
\resultmetric$=$False rows also include evaluation-side errors rather than pure
threshold misses. For these reasons, we present the paper as a measurement note,
not as a universal scaling law.

\subsection{How This Connects to Richer Agent Analysis}

The result here is deliberately simpler than full trajectory analysis. Recent
work such as \texttt{AgentRx} and \texttt{TRAJEVAL} asks why an agent failed at a
fine-grained trajectory level \cite{agentrx,trajeval}. Our claim comes earlier
in the analysis pipeline. Before manual trajectory diagnosis, the official
benchmark outputs already justify more structured reporting. That makes the
\processmetric/\resultmetric split a useful bridge: richer trajectory work can
build on it, but we should not ignore the information already provided by the
benchmark.
